\scanpage{86}\mspage{30}
% L1-4
(on a appliqué $p$), on lui applique $\Phi_0$, \deleted{\unsure{on prend}} i.e. on regarde $\Phi C X$ comme objet de $\underline{\mathrm{Fibsc}}\,S$ (et $=$ \unsure{de|dans} $\underline{\mathrm{Fibsc}}_0\,S$). On veut définir un \unsure{mor}
% L5
\[
 X\xrightarrow{\beta_X}\Phi_{\unsure{0}}p(X)=\Phi C(X)\quad\inserted{\text{dans }\text{\uuline{$\underline{\mathrm{Fibsc}}_0$}}}
\]
% L6-7
fonctoriel en $X$. Or on a comme tantôt
\[
 X\mathrel{\overset{j_X}{\unsure{\mathord{-}|\longrightarrow}}} C(X)
\]
% L8-9
d'où en appliquant $\Phi$
\[
 \Phi(X)\xrightarrow{\Phi(j_X)}\Phi C(X),
\]
% L10-15
mais \unsure{regardé|regardée} cette fois comme flèche \unsure{dans|de} $\underline{\mathrm{Fibsc}}_0\,S$ \deleted{\unsure{, d'où}} (N.B. on a bien $\Phi(X)$ dans $\underline{\mathrm{Fibsc}}_0\,S$, car $X$ \deleted{étant dans} provenant de $\underline{\mathrm{Fibsc}}_0\,S$, \deleted{\unsure{donc de}} est dans $\underline{\mathrm{Fibhot}}\,S$, et on applique la remarque ``Cor.~3'' page 21.) D'autre part, on a une flèche de $\mathrm{Cat}_{/S}$
% L16
\[
 X\xrightarrow{i_X}\Phi(X),
\]
% L17-21
est-elle une flèche dans $\underline{\mathrm{Fibsc}}_0\,S$, i.e. est-elle cartésienne ? On voit tout de suite que non, qu'on a seulement, pour une flèche $u:s\longrightarrow t$ dans $S$, un homomorphisme de commutation (pas un iso) dans
% L22-24
\diagram{\begin{tikzcd}[ampersand replacement=\&,row sep=large,column sep=large]
 X_s\arrow[r,"{u_*^{X}}"]\arrow[d,"i_s"] \& X_t\arrow[d,"i_t"{name=rightside}]\\
 X_{/s}\arrow[r,"{(\mathrm{id}_X)_{/u}}"{name=bottomside},"{\substack{\Vert\\u_*^{\Phi(X)}}}"'] \& X_{/\unsure{t|f}}
 \arrow[from=bottomside,to=rightside,bend left=12,"\alpha"]
\end{tikzcd}}
% archive sheet number: 85
